% #############################################################################
% This is Chapter 5
% !TEX root = ../main.tex
% #############################################################################
% Change the Name of the Chapter i the following line
\fancychapter{Problem Analysis}
% The following line allows to ref this chapter
\label{chap:ProbAnalysis}

The main problem this work attempts to solve is informing someone of what actually changes, operationally speaking, due to legislation. This is an important task, since a lot of people affected by legislation know nothing about how they are being affected by it. The objective is to produce artifacts that ease the understanding of how a legislative act affects that stakeholder, specifically. Stakeholder-Specific viewpoints are thus important to achieve this, given that such particular Views must be produced through the viewpoint of the stakeholder. \ac{GenAI} can be a powerful tool to assist in correctly producing these viewpoints in a timely manner.

The following chapter covers the identified limitations on the work this paper succeeds, as well as the previously identified problems. Through-out the chapter, requirements arising from the analysis of the problem will be identified. It is relevant to note the different actors that take a part in the process: The analyst or enterprise architect that utilizes the method is identified in the following text as ``user'', and the person whose viewpoint is being identified is the ``stakeholder''.

% #############################################################################
% =============================================================================
% =============================================================================
% #############################################################################
\section{Limitations on Previous Work}

The existing methodology provides a structured approach for translating legal requirements into ArchiMate architecture descriptions through traceability matrices, hierarchical identifiers, and modular views. However, important limitations affect the strength of its conclusions and its practical scalability. First, the manual construction of large-scale architecture descriptions requires substantial time and creates a residual risk of omissions, inconsistent interpretations, and modelling errors. Secondly, the methodology lacks external validation from legal experts and intended stakeholders, since its demonstrations are illustrative rather than field-tested. The following subsections examine these limitations in greater detail.
% #############################################################################
\subsection{Time and error risk in large-scale manual modelling}

The methodology faces a scalability limitation arising from the manual construction of more than one hundred views and over one thousand numbered elements. Creating, interpreting, and maintaining this volume of legal-to-ArchiMate mappings requires substantial time and sustained expert attention. As model size increases, the risk of omitted clauses, inconsistent interpretations, incorrect relationships, duplicated elements, and documentation mismatches also increases. Even when traceability identifiers are used, they do not eliminate the need to verify whether each legal provision has been represented faithfully and consistently across the full architecture description. The approach is therefore valuable for structured legal modelling, but its time cost and residual error risk may limit its practical adoption unless supported by systematic review procedures, automation, and independent quality checks. 

This work leverages on the previously explored approach and makes an attempt at partially automating it through the use of \ac{GenAI}, with pre-defined approval gates for human validation. It is also pertinent to note that the focus changes from representing an entire regulation and its intricacies to representing only what clearly affects the operational footprint of a single stakeholder. To communicate a more global vision of the regulation, in a manner similar to the previous work, an ``overview'' View is generated. The objective is to let the stakeholder know there is more to the regulation than what produces a direct impact on one's reality.

% #############################################################################
\subsection{External validation}

The previous study did not include a formal external-validation phase. The same holds true about this work, due to time constraints. Its practical scenarios were intended primarily to demonstrate how the methodology could be applied to construct traceable architecture descriptions, rather than to provide stakeholder-based validation of the resulting models.

Accordingly, the results provide evidence of the methodology's feasibility and potential usefulness, while leaving its interpretation and practical use by legal practitioners, enterprise architects, and regulated organisations open for further study. The AI-assisted reviews offered exploratory feedback during the work, but do not constitute independent legal review or stakeholder validation. External assessment of legal accuracy, usability, and real-world effectiveness therefore remain a relevant avenue for subsequent research.

\section{Stakeholder-Specific Viewpoints}

The legal and operational scope of a normative act is inherently broad. A single regulation may affect multiple organisational units, systems, processes, and governance mechanisms. An architectural representation intended to capture this scope would necessarily encompass many elements that are not relevant to the daily work or concerns of a particular stakeholder.

This creates a practical communication problem. A regulation-wide representation may be valuable from a legal or enterprise perspective, but it can be difficult for an individual stakeholder to determine which aspects affect their own responsibilities \cite{KosenkovEtAl2025}. A back-end developer may need to understand how regulatory requirements constrain code changes, testing, deployment, monitoring, and incident escalation. In contrast, a team leader may be concerned with approval responsibilities, delivery dependencies, risk ownership, and resource allocation. Presenting both stakeholders with the same broad representation risks obscuring these distinct concerns with irrelevant detail \cite{BreauxEtAl2008}.

ISO/IEC/IEEE 42010 addresses this problem through architecture viewpoints and views. A viewpoint establishes the conventions for constructing and interpreting a representation for specified stakeholders and concerns, while a view applies those conventions to present the relevant part of an architecture description \cite{ISO42010:2022}. ArchiMate similarly supports the creation of stakeholder-oriented views that select and relate the concepts needed to address a particular concern \cite{opengroup}. 

The value of a viewpoint lies in making the intended purpose of an architectural representation explicit. It identifies the stakeholders and concerns to be addressed and establishes the concepts, relationships, and presentation conventions through which those concerns are represented \cite{ISO42010:2022}. A viewpoint therefore supports the construction of views that are meaningful to a particular audience, rather than assuming that one universal representation is equally useful to all stakeholders. In order to be meaningful, they must be understood by the person they are meant to inform. This clearly states a requirement the method must fulfil, denoted as \textbf{R1:} ``Artefacts are comprehensible by the stakeholder''.

The methodology therefore constructs stakeholder-specific viewpoints around the stakeholder's documented responsibilities, dependencies, and concerns. Each viewpoint specifies how the regulatory and operational context relevant to those concerns is to be represented, and guides the creation of corresponding ArchiMate views, which support regulatory understanding and discussion.

\section{Legal to Operational Relevance Gap}

European legislation commonly formulates duties at the level of regulated entities, their management bodies, or other legally defined roles. It does not normally specify how an individual employee should perform their daily work. For example, \ac{DORA} assigns responsibilities for \ac{ICT} risk management to financial entities and their management bodies, rather than to roles within an organization \cite{DORA}. The practical effect of a legal requirement on an individual stakeholder must therefore be established through the organisational and operational context in which that stakeholder works.

The legal-to-operational relevance gap is primarily a mismatch of abstraction, terminology, and perspective. Legislation expresses normative expectations through duties, conditions, accountability, and legally defined entities or roles. Operational practice is instead described through work activities, decisions, processes, systems, information flows, and artefacts. Relating the two requires establishing how a legally expressed requirement manifests within a particular organisational context and how that context connects to a stakeholder's work. These connections are rarely stated explicitly in either the legal text or the stakeholder's operational description, which makes regulatory requirements engineering a necessary interpretative step \cite{Ghanavati2014}.

\section{Controlled Use of Generative AI}

The legal-to-operational relevance gap makes \ac{GenAI} a potentially useful analytical aid. The process requires the examination of extensive legal text, the structuring of stakeholder-provided operational information, and the comparison of concepts expressed in substantially different language. \ac{GenAI} can assist with these labour-intensive activities by extracting candidate facts, organising information, identifying potential relationships, and producing structured representations for review. Such uses align with emerging research on AI-supported regulatory requirements analysis \cite{Abualhaija2024}.

However, generated output cannot itself be regarded as legal, operational, or architectural evidence. \ac{GenAI} systems may produce plausible but erroneous statements, including content that deviates from supplied sources or presents unsupported conclusions with confidence. \ac{NIST} identifies this risk as confabulation, commonly referred to as hallucination \cite{NISTGenAI2024}. In the present context, this could result in a fabricated legal reference, an unsupported claim that a requirement applies to a stakeholder, or an architectural relationship that is based on general domain knowledge rather than documented organisational evidence. These can be mitigated by implementing mandatory human approval-gates, meant to validate the produced output (\textbf{R2:}``Human approval required'').

To add to this issue, \ac{LLM}s are also non-deterministic. The same input will never produce the same exact output twice. In the scope of the current work, this can be an issue when identifying what requirements enforced by a regulation apply to the operational footprint of a stakeholder, specially when dealing with a requirement that applies in a more indirect manner. These eventual divergences can be mitigated through the use of an agent debate panel, previously discussed in \ref{sec:multi-agent-debate} and identify a new requirement \textbf{R3:} ``Produced results must be stable''.

The use of \ac{GenAI} must therefore be controlled through explicit methodological boundaries. Legal statements must remain traceable to a verified legal source (\textbf{R4:} ``Traceability to the legal source''), while claims regarding stakeholder responsibilities, activities, and dependencies must remain traceable to documented operational evidence (\textbf{R5:}``Traceability to operational evidence'') . Where the available sources do not substantiate a relationship, the methodology must not allow the system to infer one ( \textbf{R6:} Nothing is inferred). \ac{GenAI} should support the identification and structuring of candidate information, rather than independently determining legal applicability, operational relevance, or organisational compliance.

Human involvement is consequently a substantive validation activity rather than a final confirmation step. A human reviewer must be able to inspect the evidence and rationale behind each material conclusion, reject unsupported output, and explicitly approve the information used in subsequent stages. This preserves accountability for interpretative decisions while allowing \ac{GenAI} to reduce the effort required to produce evidence-traceable stakeholder viewpoints. The methodology therefore employs \ac{GenAI} as a controlled analytical aid, not as an autonomous compliance assessor.